%------------------------------------------------------------------------------%
\documentclass[fleqn,utf8,aspectratio=169,12pt]{beamer}

\usepackage{calculo_varias_variaveis-1}

\title{Vetor Tangente}

%------------------------------------------------------------------------------%
\begin{document}

\addtitle

%------------------------------------------------------------------------------%
\section{Vetor}

%------------------------------------------------------------------------------%
\begin{frame}{Vetor}
\begin{itemize}[<+->]
  \setlength\itemsep{1em}
  \item Vetor é definido por \emph{módulo}, \emph{direção} e \emph{sentido}
  \item Representa uma classe de \emph{segmentos de reta orientados} equipolentes\\
        (paralelos, mesmo tamanho e sentido)
  \item É representado por componentes \((x,y)\) ou \((x,y,z)\)
  \item Pode ser posicionado em qualquer ponto do plano ou espaço
\end{itemize}
\end{frame}

\include{ex/B-03-Vetor_tangente-ex0}

%------------------------------------------------------------------------------%
\section{Vetor Tangente}

%------------------------------------------------------------------------------%
\begin{frame}{Curva Paramétrica}
  $x$ e $y$ são dados como funções
  \[
    x = f(t)     \qquad
    y = g(t)     \qquad
    t \in I
  \]

  \pause
  Vetorialmente
  \[
    \gamma(t)
    \;=\; \vetor{x}{y}
    \;=\; \vetor{f(t)}{g(t)}
  \]
\end{frame}

\framedgraphic{Vetor Tangente}{B_03_vetor_tangente}

%------------------------------------------------------------------------------%
\begin{frame}{Vetor Tangente}
  Se
  \[
    \gamma(t)
    = \vetor{x}{y}
    = \vetor{f(t)}{g(t)}
    \qquad
    t \in I
  \]
  \pause
  e $f$ e $g$ forem deriváveis em $t$,
  \pause
  o \emph{vetor tangente} é
  \pause
  \[
    \gamma'(t)
    \pause
    \;=\; \frac{d}{dt} \vetor{x}{y}
    \pause
    \;=\; \vetor{x}{y}'
    \pause
    \;=\; \Vetor{\dfrac{df}{dt}}{\dfrac{dg}{dt}}
    \pause
    \;=\; \vetor{f'(t)}{g'(t)}
  \]
\end{frame}

%------------------------------------------------------------------------------%
\section{Exemplos}

%------------------------------------------------------------------------------%
\begin{question}
\begin{columns}
\column{0.5\linewidth}
  Calcule o vetor tangente à curva
  \[
    x = t - t^2 \qquad
    y = t - t^3
  \]
\column{0.5\linewidth}
  \pause
  ~\\[5mm]
  \includegraphics{B_03_exemplo_1_curva}
\end{columns}
\end{question}

%------------------------------------------------------------------------------%
\begin{resolution}{Solução}
\begin{columns}
\column{0.5\linewidth}
  Calculando \(\frac{dx}{dt}\) e \(\frac{dy}{dt}\)

  \pause
  \medskip
  \[
    \frac{dx}{dt}
    \pause
    = \frac{d}{dt}\left(t - t^2\right)
    \pause
    = 1 - 2t
  \]

  \pause
  \medskip
  \[
    \frac{dy}{dt}
    = \frac{d}{dt}\left(t - t^3\right)
    \pause
    = 1 - 3t^2
  \]
\column{0.5\linewidth}
  \pause
  Vetor tangente
  \pause
  \medskip
  \[
    v(t)
    \pause
    = \Vetor{\frac{dx}{dt}}{\frac{dy}{dt}}
    \pause
    = \vetor{1 - 2t}{1 - 3t^2}
  \]
\end{columns}
\end{resolution}

%------------------------------------------------------------------------------%

\include{ex/B-03-Vetor_tangente-ex2}
%------------------------------------------------------------------------------%
\begin{question}
\begin{columns}
\column{0.5\linewidth}
  Calcule o vetor tangente à curva
  \[
    x = t\cos(t)
  \]

  \[
    y = t\sin(t)
  \]

  \[
    0 \leq t < \infty
  \]
  quando \(t=\frac{\pi}{2}\)
\column{0.5\linewidth}
  \pause
  ~\\[5mm]
  \includegraphics{B_03_exemplo_3_curva}
\end{columns}
\end{question}

%------------------------------------------------------------------------------%
\begin{resolution}{Solução}
  \begin{columns}
    \column{0.5\linewidth}
    \onslide<+->{Derivada da coordenada $x$}
    \begin{align*}
      \onslide<+->{x'(t) & = \frac{dx}{dt}                \\}
      \onslide<+->{      & = \frac{d}{dt}t\cos(t)         \\}
      \onslide<+->{      & = 1\times\cos(t) + t(-\sin(t)) \\}
      \onslide<+->{      & = \cos(t) - t\sin(t)             }
    \end{align*}

    \column{0.5\linewidth}
    \onslide<+->{Derivada da coordenada $y$}
    \begin{align*}
      \onslide<+->{y'(t) & = \frac{dy}{dt}              \\}
      \onslide<+->{      & = \frac{d}{dt}t\sin(t)       \\}
      \onslide<+->{      & = 1\times \sin(t) + t\cos(t) \\}
      \onslide<+->{      & = \sin(t) + t\cos(t)           }
    \end{align*}
  \end{columns}
\end{resolution}

%------------------------------------------------------------------------------%
\begin{resolution}{Solução}
  \onslide<+->{Avaliando as funções derivada em \(t=\frac{\pi}{2}\)}
  \[
    \onslide<+->{x'\left(\frac{\pi}{2}\right)}
    \onslide<+->{   = \left(\cos(t) - t\sin(t)\right) \evalat{\frac{\pi}{2}}{}     }
    \onslide<+->{   = \cos\left(\frac{\pi}{2}\right) - \frac{\pi}{2}\sin\left(\frac{\pi}{2}\right)}
    \onslide<+->{   = 0 - \frac{\pi}{2} \times 1}
    \onslide<+->{   = -\frac{\pi}{2}}
  \]
  \[
    \onslide<+->{y'\left(\frac{\pi}{2}\right)}
    \onslide<+->{   = \left(\sin(t) + t\cos(t)\right) \evalat{\frac{\pi}{2}}{}     }
    \onslide<+->{   = \sin\left(\frac{\pi}{2}\right) + \frac{\pi}{2}\cos\left(\frac{\pi}{2}\right)}
    \onslide<+->{   = 1 + \frac{\pi}{2} \times 0}
    \onslide<+->{   = 1}
  \]

  \bigskip
  \onslide<+->{O vetor tangente é
    \(
    \begin{bmatrix}
      \sfrac{-\pi}{2} \\ 1
    \end{bmatrix}
    \)}
\end{resolution}

%------------------------------------------------------------------------------%
\begin{resolution}{Solução}
  \centering
  \includegraphics{B_03_exemplo_3_vetor}
\end{resolution}

%------------------------------------------------------------------------------%
%------------------------------------------------------------------------------%
\begin{question}
\begin{columns}
\column{0.5\linewidth}
  Calcule o vetor tangente à curva
  \[
  \begin{cases}
    x & \!\!\!\!= e^{2t} \cos(\pi t) \\
    y & \!\!\!\!= e^t \sin(\pi t)
  \end{cases}
  \]
  em \(t=\frac{3}{4}\)
\column{0.5\linewidth}
  \pause
  ~\\[5mm]
  \includegraphics{B_03_exemplo_4_curva}
\end{columns}
\end{question}

%------------------------------------------------------------------------------%
\begin{resolution}{Solução}
  \onslide<+->{Derivada da coordenada $x$}
  \begin{align*}
    \onslide<+->{x'(t) & = \frac{dx}{dt}       \\}
    \onslide<+->{      & = \frac{d}{dt}\left(e^{2t} \cos(\pi t)\right) } \\
    \onslide<+->{      & = \frac{d}{dt}\left(e^{2t}\right) \cos(\pi t)
      + e^{2t} \frac{d}{dt}\left(\cos(\pi t)\right) } \\
    \onslide<+->{      & = e^{2t}\frac{d}{dt}\left(2t\right) \cos(\pi t)
      - e^{2t} \sin(\pi t)\frac{d}{dt}\left(\pi t\right) } \\
    \onslide<+->{      & = 2 e^{2t} \cos(\pi t) - \pi e^{2t} \sin(\pi t) }
  \end{align*}
\end{resolution}

%------------------------------------------------------------------------------%
\begin{resolution}{Solução}
  \onslide<+->{Derivada da coordenada $y$}
  \begin{align*}
    \onslide<+->{y'(t) & = \frac{dy}{dt}       \\[2mm]}
    \onslide<+->{      & = \frac{d}{dt}\left(e^t \sin(\pi t)\right)} \\
    \onslide<+->{      & = \frac{d}{dt}\left(e^t\right) \sin(\pi t)
      + e^t \frac{d}{dt}\left(\sin(\pi t)\right)} \\
    \onslide<+->{      & = e^t \sin(\pi t)
      + e^t \cos(\pi t)\frac{d}{dt}\left(\pi t\right)} \\
    \onslide<+->{      & = e^t \sin(\pi t)+ \pi e^t \cos(\pi t)}
  \end{align*}
\end{resolution}

%------------------------------------------------------------------------------%
\begin{resolution}{Solução}
  \onslide<+->{Avaliando a derivada de $x(t)$ em \(t=\sfrac{3}{4}\)}
  \begin{align*}
    \onslide<+->{x'\left(\frac{3}{4}\right)}
    \onslide<+->{& = \left[ 2 e^{2t} \cos(\pi t) - \pi e^{2t} \sin(\pi t)\right]
                     \evalat{t=\sfrac{3}{4}}{}}
    \\
    \onslide<+->{& =   2 e^{2\frac{3}{4}} \cos\left(\pi \frac{3}{4}\right)
                   - \pi e^{2\frac{3}{4}} \sin\left(\pi \frac{3}{4}\right)}
    \\
    \onslide<+->{& =   2 e^{\sfrac{3}{2}} \cos\left(\frac{3\pi}{4}\right)
                   - \pi e^{\sfrac{3}{2}} \sin\left(\frac{3\pi}{4}\right)}
    \\
    \onslide<+->{& =   2 e^{\sfrac{3}{2}} \left(-\frac{\sqrt{2}}{2}\right)
                   - \pi e^{\sfrac{3}{2}} \frac{\sqrt{2}}{2}}
    \\
    \onslide<+->{& = - \sqrt{2} \left(1 + \frac{\pi}{2} \right) e^{\sfrac{3}{2}}}
  \end{align*}
\end{resolution}

%------------------------------------------------------------------------------%
\begin{resolution}{Solução}
  \onslide<+->{Avaliando a derivada de $y(t)$ em \(t=\sfrac{3}{4}\)}
  \begin{align*}
    \onslide<+->{y'\left(\frac{3}{4}\right)}
    \onslide<+->{& = \big(e^t \sin(\pi t)+ \pi e^t \cos(\pi t)\big)
      \evalat{t=\sfrac{3}{4}}{}}
    \\
    \onslide<+->{& =     e^{\sfrac{3}{4}} \sin\left(\frac{3\pi}{4}\right)
                   + \pi e^{\sfrac{3}{4}} \cos\left(\frac{3\pi}{4}\right)}
    \\
    \onslide<+->{& =     \frac{e^{\sfrac{3}{4}}}{\sqrt{2}}
                   - \pi \frac{e^{\sfrac{3}{4}}}{\sqrt{2}}}
    \\
    \onslide<+->{& = \frac{1 - \pi}{\sqrt{2}} e^{\sfrac{3}{4}}}
  \end{align*}
\end{resolution}

%------------------------------------------------------------------------------%
\begin{resolution}{Solução}
  O vetor tangente no ponto \(t=\sfrac{3}{4}\) é
  \[
    \frac{d}{dt}\left(\!
    \begin{array}{c}
      x \\ y
    \end{array}
    \!\right)
    =
    \left(
    \begin{array}{c}
      - \sqrt{2} \left(1 + \dfrac{\pi}{2} \right) e^{\sfrac{3}{2}} \\[9mm]
        \dfrac{1 - \pi}{\sqrt{2}} e^{\sfrac{3}{4}}
    \end{array}
    \right)
  \]
\end{resolution}

%------------------------------------------------------------------------------%
\begin{resolution}{Solução}
  \centering
  \includegraphics{B_03_exemplo_4_vetor}
\end{resolution}

%------------------------------------------------------------------------------%

\include{ex/B-03-Vetor_tangente-ex5}

%------------------------------------------------------------------------------%
\section{Lista Mínima}

%------------------------------------------------------------------------------%
\begin{frame}{Lista Mínima}
  Cálculo Vol. 2 do Thomas 12\textsuperscript{a} ed. -- Seção 11.2
  \begin{enumerate}
    \item Estudar todo o texto da seção
    \item \emph{Calcule o vetor tangente}
          das funções dos exercícios: 2, 8, 12, 16, 18, 20\\
          (Ignore o enunciado original dos exercícios)
  \end{enumerate}

  \vfill
  Atenção: A prova é baseada no livro, não nas apresentações
\end{frame}

%------------------------------------------------------------------------------%
\end{document}
%------------------------------------------------------------------------------%
